\chapter{原子观念是怎样被逼出来的}

\begin{kaichang}
去厨房拿一块方糖（没有方糖，一小颗冰糖也行）。现在做一件看起来很无聊的事：把它掰成两半，再把其中一半掰成两半，再掰，再掰……

用不了几次，糖粒就小到捏不住了。假设你有无限的耐心和一双无限灵巧的手，还有一个永远能放大细节的显微镜——请问：这条路能一直走下去吗？糖能不能被无限地分下去，永远分得出来“更小的糖”？

先别往下看。在纸上写下你的答案：能，或者不能。这个两千多年前就有人认真吵过的问题，最后不是靠哲学家吵架吵赢的，而是被天平、量筒和一堆冷冰冰的数字逼出了答案。这一章就讲这个过程。
\end{kaichang}

\section{一个吵了两千年的问题}

古希腊的哲学家德谟克利特猜想：万物都由极小、极硬、不可分割的颗粒组成，他把这种颗粒叫做“原子”（希腊语原意就是“切不开”）。同一片土地上，亚里士多德一派却认为物质是连续的，可以无限细分，万物不过是土、水、气、火四种“性质”的组合。

谁对？在当时，这是一场没法裁决的辩论。德谟克利特的猜想听起来更接近今天的答案，但请注意：\textbf{接近答案不等于科学}。他没有拿出任何一条能让怀疑者动手检验的证据，也没有说出“如果原子存在，我们应该观察到什么、观察不到什么”。这样的猜想可以很漂亮，却没办法被证实或被推翻——而“能被检验”恰恰是科学和玄想的分界线。原子观念要真正成为科学，必须等到有人把它逼到墙角：要么解释实验数字，要么出局。

被逼问的时刻，在十八、十九世纪之交到来了。逼问它的工具，是你家里也有的东西——秤。

\section{天平上的三条“死规律”}

十八世纪的化学家（拉瓦锡是代表）开始带着一种“会计心态”做实验：反应前后，每样东西都称一称。化学变化于是从“看热闹”变成了“对账”。账一对，三条规律浮出水面。

\textbf{第一条：质量守恒。}把化学反应在密闭容器里进行，反应前后总质量分毫不差。烧一根蜡烛，蜡好像“消失了”，可如果把生成的二氧化碳和水蒸气全部收集起来一起称，账是平的。物质不会凭空来，也不会凭空走。

\textbf{第二条：定比定律。}不管水是从河里取的还是实验室里制的，把它分解，得到的氢气和氧气的质量比永远是 $1:8$。你在家也能验证它的“顽固”：烧水时壶里留下白色水垢，可壶里剩下的水，以及你烧开又冷凝收集的蒸馏水，分解起来照样是 $1:8$——杂质可以浑水摸鱼，水本身的配方却雷打不动。氧化铜呢？随便你怎么制备，里面铜和氧的质量比总是约 $4:1$。一种纯净的化合物，组成比例雷打不动。

\textbf{第三条：倍比定律。}碳和氧能组成两种东西：一氧化碳（有毒的那个）和二氧化碳。取同样多的碳，比如都是 \ud{12}{g}，前者配上 \ud{16}{g} 氧，后者配上 \ud{32}{g} 氧——\textbf{恰好是 $1:2$}，不多不少，干干净净的整数比。氮和氧组成的几种化合物里，氧的比例依次是 $1:2:4$……全是整数。

请你停一秒，感受一下第三条有多古怪。如果物质像面团一样是连续的、可以任意分割的，那么同样的碳凭什么“偏偏”配上整数倍的氧？它完全可以配上 $1.37$ 倍、$2.61$ 倍。自然界却对整数情有独钟。这就像你去买米，发现米店只按整粒卖，绝不卖半粒——你会立刻意识到：\textbf{米是一粒一粒的}。

\section{道尔顿：把小颗粒请回科学}

1803 年前后，英国教师道尔顿正是这样想的。他把三条定律摆在一起，发现只要做一个假设，它们全部自动成立：

\begin{itemize}[itemsep=2pt]
\item 每种元素都由一种极小的、在化学变化中不可分割的颗粒（原子）组成；
\item 同种元素的原子，质量和性质完全相同；不同元素的原子，质量不同；
\item 化合物由不同种的原子按\textbf{简单整数比}结合而成；化学反应只是原子重新排队，原子本身不多不少、不增不灭。
\end{itemize}

现在回头看三条定律，它们全成了推论：原子在反应中不增不灭，质量当然守恒；水分子总是“1 份氢原子团配 1 份氧原子团”，比例当然固定；一个碳原子配一个氧原子，或者配两个氧原子，氧当然是 $1:2$——原子只能整颗整颗地搭配，就像米只能整粒卖。

注意这里发生的方法论转折。德谟克利特的原子是一句“我觉得”，道尔顿的原子是一个\textbf{能算账的模型}：它不仅解释已有数据，还能预言——如果倍比定律在某个化合物上失败，原子论就得出局。一个理论肯把自己的死法写在脸上，它才是科学。

当然，道尔顿的账也不是一开始就全对。他手头没有分子式的概念，曾以为水是“一个氢原子配一个氧原子”，写出的式子相当于今天的 $\chem{HO}$；他还把氨写成 $\chem{NH}$。错在哪里？他只知道质量比，不知道每种原子一颗有多重。这个漏洞后来由意大利人阿伏伽德罗的假说补上：同温同压下，同体积的气体含有同样多的粒子。对比反应前后气体的体积，就能反推出“水里氢氧是 $2:1$ 而不是 $1:1$”。你看，连原子论的奠基人也在自己搭的梯子上踩空过——模型不是一次到位的，是被数据一次一次修正的。

\begin{qiaoliang}
原子论一出场就带着一个让人头皮发麻的推论：原子的个数多到离谱。来算一笔账（第 5 章已经认识过“摩尔”这个计数单位）。一小滴水大约 \ud{0.05}{g}，里面有多少个水分子？
\[
\frac{\ud{0.05}{g}}{\ud{18}{g/mol}}\times \ud{6.02\times 10^{23}}{mol^{-1}} \approx 1.7\times 10^{21}\ \text{个}
\]
一千七百亿亿个——写全了是 17 后面 20 个零。这个数字大到什么程度？把全世界海洋、湖泊、河流里的水，一杯一杯（每杯 \ud{200}{mL}）分装出来，总杯数大约是 $10^{21}$ 量级。也就是说，\textbf{一小滴水里的水分子数，比地球上所有水能分出的杯数还多}。反过来看，每个水分子小到什么程度，你自己除一下就知道了。
\end{qiaoliang}

\section{符号姗姗来迟}

有了“原子按整数比搭配”的图景，记账符号就水到渠成了。道尔顿自己画过一套圆圈符号——氧是一个空心圆，氢是带点的圆，水是两个圆挨在一起——繁琐得像画连环画。后来瑞典化学家贝采利乌斯提议：干脆用拉丁名称的首字母。O 代表氧的一个原子，H 代表氢的一个原子，C 代表碳的一个原子。

于是 $\chem{H_2O}$ 这串符号的真正含义就清楚了：\textbf{这不是缩写，是人口普查}——它报告“一个水粒子由 2 个氢原子和 1 个氧原子组成”。$\chem{CO}$ 和 $\chem{CO_2}$ 的区别也不在于写法，而在于碳原子搭配氧原子的颗数不同：$1:1$ 对 $1:2$，正是倍比定律里那个整数比。符号是粒子的影子，先有粒子，后有符号——所以这本书一直坚持到最后才写符号。

\section{还有人就是不信}

道尔顿的原子论解释了一切，可整个十九世纪，始终有一批重量级的化学家和物理学家拒绝买账，比如德国化学家奥斯特瓦尔德。他们的理由并非胡搅蛮缠：原子看不见、摸不着，谁知道是不是一套碰巧好使的记账技巧？化学反应的整数比，也许只是某种连续规律的巧合。

这场争论持续到二十世纪初。终结它的，是显微镜下一粒花粉的“舞蹈”。

\begin{zhengju}
1827 年，植物学家布朗在显微镜下观察悬浮在水中的花粉颗粒，发现它们一刻不停地做无规则抖动，像被看不见的手推来搡去。他换用无生命的石粉，抖动依旧——排除了“花粉是活的”这个解释。此后的七十多年里，没人能说清这股推力从哪来。

1905 年，爱因斯坦给出了一个大胆的解释：花粉颗粒被周围的水分子从四面八方撞击。分子太小看不见，但亿万个分子的撞击不可能每时每刻都恰好抵消，总有一瞬间某一侧撞得更狠，颗粒就被“推”了一下。他更进一步算出：如果水真的由这种大小的分子组成，花粉颗粒在单位时间内应该“漂移”多远——一个可以拿到显微镜下去量的具体数字。

法国物理学家佩兰真的去量了。他花了好几年，用显微镜追踪悬浮颗粒的轨迹，测出的漂移幅度与爱因斯坦的预言吻合，还由此独立算出了每摩尔粒子的数目。账对上了。1908 年之后，连奥斯特瓦尔德都公开认输，承认原子和分子是真实存在的。这个故事值得记住的不是某个天才，而是这个链条：一个古怪的观测（花粉抖动）$\rightarrow$ 一个可检验的定量解释（分子撞击）$\rightarrow$ 一场精心设计、长达数年的测量（佩兰）$\rightarrow$ 最顽固的怀疑者改口。科学共识就是这样炼成的，不靠权威点头。
\end{zhengju}

\begin{shiyan}
\textbf{在家观察“分子的撞击”}

材料：一杯清水、一滴稀释的红墨水（或碾碎的一点水彩笔芯）、白纸、放大镜（或手机微距模式）。

步骤：把一滴墨水轻轻滴进静置的清水表面，不要搅动，把杯子放在白纸上观察十分钟。

观察点：墨迹不是均匀地摊开，而是伸出丝丝缕缕、弯弯曲曲的“触手”，方向毫无规律。这就是水分子无规则运动推搡染料颗粒留下的痕迹。

进阶：接一杯热水、一杯冷水，各滴一滴墨水同时对比——热水里扩散明显更快。温度越高，分子运动越猛烈，这正是“温度是分子运动剧烈程度”的直接目击（第 4 章还会回来用它）。

安全提示：全程只用清水和文具，无须加热明火，结束后正常清洗即可。
\end{shiyan}

\section{“切不开”的粒子被切开了}

道尔顿假设原子不可分割，这个假设在化学实验里从没失手。但十九世纪末，几样新玩具——真空管、高压电、感光底片——把原子内部捅开了三个口子。

\textbf{第一个口子：电子（1897 年）。}汤姆孙研究阴极射线管——一根抽了真空的玻璃管，两端加高压电，阴极会发出一种看不见的“射线”，能在玻璃上打出荧光。他让这束射线穿过电场和磁场，发现它总是朝带正电的板偏转，而且偏转程度与管子里装什么气体、用什么金属做电极\textbf{毫无关系}。结论很硬：这是一种带负电的微粒，它比最轻的原子（氢）还要轻上千倍，并且是\textbf{所有原子共有的零件}。原子里面居然有更小的东西——“切不开”这三个字，从此打了折扣。

\textbf{第二个口子：原子核（1911 年）。}汤姆孙的原子画像很自然：原子是一团均匀的正电荷“布丁”，电子像葡萄干一样嵌在里面。他的学生卢瑟福决定检验一下。他让一束带正电的高速粒子（\chem{\alpha} 粒子，后来被认出是氦原子核）轰击极薄的金箔。如果原子真是均匀的布丁，这些子弹般的粒子应该几乎不受影响地穿过去——绝大多数确实如此。但大约每两万个粒子里，有一个被\textbf{猛地弹了回来}，偏转角度大得离谱。卢瑟福后来说，这就像朝一张纸巾开炮，结果炮弹被弹回来砸在你脸上。

能让高速粒子弹回来，只有一种可能：原子的几乎全部质量、以及全部正电荷，都挤在一个极小的核心里。原子不是布丁，而是“空空荡荡的房子”：中央一个极小极重的核，电子在核外广大空间里活动。这个结论之所以让人信服，还因为它给出了一串可检验的定量预言——偏转角越大，被弹回的粒子应该越少，具体少多少可以算出来。卢瑟福的助手盖革和学生马斯登数了几百万次闪光，数字与预言逐点吻合。

\textbf{第三个口子：质子和中子（1919、1932 年）。}卢瑟福后来用 \chem{\alpha} 粒子轰击氮气，打出了氢核——也就是质子，原子核的基本成员之一。但质子数对不上核的质量：比如氦核有 2 个质子，质量却接近 4 个质子。谜底由查德威克在 1932 年揭开：核里还有一种不带电、质量与质子几乎相同的粒子——中子。至此，原子的“全家福”凑齐：质子、中子挤在核里，电子住在核外。至于质子数怎样决定“你是谁”（金还是铁），那是第 7 章“元素身份证”的故事。

\section{三个模型，三次胜利，三次露馅}

值得停下来看一眼这条模型接力赛，因为每一棒都漂亮，也都留了破绽：

\begin{itemize}[itemsep=2pt]
\item \textbf{汤姆孙模型（葡萄干布丁）}：解释了原子呈电中性（正负相抵）、电子可以被“揪”出来。错在：正电荷并非均匀分布——金箔实验直接把它枪毙了。
\item \textbf{卢瑟福模型（核式模型）}：完美解释了散射实验，发现了原子核。错在：按当时的电磁学，绕核运动的电子会不断辐射能量、在极短的时间内螺旋坠进核里——原子根本活不过一眨眼。可原子明明很稳定。核式模型解释了“核存在”，却解释不了“原子为什么能活着”。
\item \textbf{玻尔模型（1913 年）}：玻尔强行规定：电子只能待在某些特定的“允许轨道”上，在这些轨道上\textbf{不}辐射能量；只有从一条轨道跳到另一条时，才以发光的方式交出能量差额。这一下同时救活了原子的稳定性，还顺手解释了氢原子发光为什么是几条特定颜色的亮线（光谱的具体证据，第 9 章会细讲）。错在：这套“规定”本质上是打补丁——为什么轨道偏偏是这些？模型答不上来；而且一遇到比氢复杂的原子，它的计算就开始失灵。
\end{itemize}

看清楚了吗？每个模型都不是“错了被扔掉”，而是\textbf{解释了一批现象，又被下一批现象逼到墙角}。今天更准确的图景——电子不是沿固定轨道跑的小行星，而是一团按概率分布的“云”——是第 9 章的主角。你此刻只需要记住接力赛的节奏：现象逼出模型，新现象再逼出新模型。原子观念不是某个人灵机一动发明出来的，是被证据一轮一轮逼出来的。

\section{原子有多小，核有多空}

把尺度数字摆在一起，你会重新理解“空旷”两个字。原子的直径大约是 $10^{-10}$ 米量级（不同元素略有差别），原子核的直径只有约 $10^{-15}$ 到 $10^{-14}$ 米——\textbf{原子核只有原子的十万分之一}。

换个人能把握的比喻：把一个原子放大到一座标准足球场那么大，原子核就是摆在球场正中的一颗玻璃弹珠，而整个“球场”的其余部分，只有若有若无的电子在活动。你此刻坐着感觉坚实无比的椅子，它的组成物质里超过 $99.99999999999\%$ 是空的。你之所以没有穿过椅子掉下去，靠的不是“填满”，而是电子之间的电磁排斥——那是第 17 章要讲的故事。

这也顺带解释了金箔实验的数字：原子核在原子横截面上占的面积比例大约是 $(10^{-14}/10^{-10})^2 = 10^{-8}$，即使金箔里叠着上千层原子，一个 \chem{\alpha} 粒子恰好“正中”某个核的概率也只有约万分之一量级——和实验里“两万分之一被大角度弹回”对得上号。数字能互相咬合，模型才值得信任。

\begin{wuqu}
“原子是实心的小球，像一颗微型弹珠。”——这是对原子最常见的脑内画面，教科书插图要负一半责任。真实情况是：原子的质量挤在体积分量不到万亿分之一的核里，核外没有“球壳”，只有按概率分布的电子。实验里也没有任何一种“看”原子表面的方法——原子的“边缘”本来就是模糊的。把原子想成弹珠，你就无法理解化学键（第 17 章）为什么能形成、金箔实验为什么大多数粒子一穿而过。插图里的实心小球只是画图的人不得不画的妥协。
\end{wuqu}

\section{这些模型的边界在哪里}

老规矩，说清楚本章的模型什么时候好用、什么时候会失灵：

\begin{itemize}[itemsep=2pt]
\item \textbf{“原子不可分割”}：在一切化学反应中都成立——烧掉、锈蚀、消化，原子只是重新排队。但在核反应里不成立：原子核会裂变、聚变、衰变（第 8 章专门讲），质子和中子本身还由更基本的夸克组成。“切不开”是化学尺度上的说法，不是绝对真理。
\item \textbf{玻尔模型}：对氢原子的光谱算得出奇地准，所以至今仍是入门的好帮手；但对多电子原子定量失效，轨道概念本身也被量子力学取代。它是梯子，不是终点。
\item \textbf{“整数比”}：定比、倍比定律对绝大多数化合物成立，但少数固体（某些金属氧化物）的组成可以在小范围内连续变化，不是严格的整数比——原子论仍然正确，只是这些晶体的“配方”允许少量原子缺席，细节属于更高阶段的课程。
\end{itemize}

\section{再看那块方糖}

回到开场的问题：糖能不能无限分下去？现在你有了完整的回答链条。糖由碳、氢、氧原子按固定比例组成（蔗糖的符号 $\chem{C_{12}H_{22}O_{11}}$ 就是一份人口普查：每颗糖分子 12 个碳原子、22 个氢原子、11 个氧原子）。一直掰下去，糖粒越来越小，但只要它还叫“糖”，最小单元就是那 45 个原子搭成的一颗糖分子；再分，就不是“更小的糖”，而是拆散原子——糖没了，化学性质归零。道尔顿的整数比、佩兰测出的粒子数，都在替这个答案担保。

生活里到处能见到这条链子的影子。铁钉生锈后重量会\textbf{增加}——不是铁“变多了”，而是氧原子从空气里赶来入伙，称一称便知账是平的；切开的苹果放久了变褐，是果肉里的分子和氧原子重新排队的结果；而把糖烧焦得到的黑色残渣留下、白色甜味消失，正是糖分子被拆散、原子各奔东西的现场。这些看得见的变化，全都服从同一本原子账簿。

而德谟克利特猜对的那一半——存在不可分割的单元——要加上三条修正：单元不是“最小的”，原子内部还有质子、中子和电子；单元不是“实心硬球”，而是空旷的结构；“不可分割”只在化学反应的舞台上成立。

\section{这栋房子我们才看到门厅}

原子观念早已冲出教科书，变成实打实的产业。1981 年发明的扫描隧道显微镜（STM），用一根尖到只剩一个原子的探针贴着金属表面扫描，通过探针与表面之间微弱的电流变化，能“摸”出表面上一颗一颗原子的排布——1990 年，IBM 的科学家用它在镍表面一颗颗搬动氙原子，拼出了三个字母的公司标志。今天，芯片工厂的工程师就是在这种原子尺度的图景里设计晶体管的：你的手机能在指甲盖大的芯片里塞进上百亿个开关，靠的正是“原子是一颗一颗、可以一颗颗安排”这个两百年前被天平逼出来的观念。

不过，这一章只回答了“原子存在吗”“原子里面大概什么样”。可门厅之后的问题更迷人：核里有几个质子，凭什么就决定了它是金还是铁、是惰性气体还是暴脾气？那张挂在化学教室墙上、一百多个格子的表，其实就是按这条线索排的（第 10 章会给你地图的全貌）。下一章，我们先去取那张决定元素身份的“身份证”——原子核里藏着的那个数字。

\begin{tiaozhan}
卢瑟福散射可以用概率直觉粗算。把金箔想成 $N$ 层原子叠起来的“靶墙”，每层里一个 \chem{\alpha} 粒子命中核的概率约等于核横截面积占原子横截面积的比例：$(10^{-14}/10^{-10})^2 = 10^{-8}$。卢瑟福用的金箔约几百到上千层原子，取 $N \approx 10^3$，则单个粒子命中核的概率约 $10^3 \times 10^{-8} = 10^{-5}$，即十万分之一量级。实验观测到的大角度散射比例约为两万分之一，考虑到“命中”也不一定弹回正后方，两者吻合得相当不错。这说明：不需要看见原子，只要“数子弹、量反弹”，就能反推出靶子内部的结构——这是整个微观物理学的看家本领。跳过本框不影响主线。
\end{tiaozhan}

\section*{练一练}

\begin{sikaoti}
\item 用道尔顿原子论解释：为什么“同样的碳，分别与 \ud{16}{g} 和 \ud{32}{g} 氧结合”这条观测，会成为原子存在的强证据？如果物质是连续的，这个比例应该长什么样？
\item 汤姆孙发现：无论真空管里装哪种气体、用哪种金属做电极，阴极射线粒子的偏转都一样。为什么“与材料无关”这一点，是“电子是所有原子共有零件”的关键证据？
\item 画一幅卢瑟福散射的示意图：画出一层金原子（大圆圈里一个小点代表核），再画五条 \chem{\alpha} 粒子的入射路线，分别让四条几乎直行、一条大角度偏折。图旁注明：圆圈大小不代表原子真实“外壳”，只是帮助思考的表示法。
\item 如果把一个氢原子放大到直径 \ud{100}{m}（约一个操场），按比例它的原子核直径大约是多少毫米？算完再想想：为什么这个比例能解释金箔实验里“绝大多数粒子直接穿过”？
\item 玻尔模型规定电子只能在特定轨道上、且待在上面时不发光。这条规定“救活”了原子的什么性质？它又为什么说不上是自然界的真正解释？
\item 有人说：“原子是古希腊人想出来的，所以科学只是运气好蒙对了。”请用德谟克利特猜想与道尔顿原子论的区别，写三句话反驳他。
\end{sikaoti}
